\section*{Part VI: The Models Are Different Because the Economies Are Different}

\subsection*{54. FMCG and Retail: Low Unit Value, High Repetition, Channel Economics}

For a consumer packaged good, the unit economics may begin with retailer selling price but must end at the manufacturer's realized net revenue after trade margins, promotions, returns, and other adjustments. The direct costs include inputs, processing, packaging, fulfillment, and waste. Replenishment, shelf placement, logistics, and repeat purchase often matter more than dramatic 1-time conversion.

Suppose a pack sells to the consumer for ₹100, but the manufacturer's net realization is ₹60 after channel economics, and production plus distribution variable expense is ₹42. The manufacturer has ₹18 of simplified contribution before promotion and fixed costs. Reporting ₹100 as manufacturer revenue or calling the ₹58 difference between shelf price and input cost ``gross margin'' would be wrong. The entrepreneur must learn which price belongs to which participant in the value chain.

Working capital is often central: inventory is purchased before sales; retailers may demand credit; spoilage and expiry can destroy realized margin. Revenue growth with expanding receivables and inventory may be cash-negative. A healthy FMCG model therefore requires both repeat demand and efficient cash conversion.

\subsection*{55. E-Commerce: The Return and Fulfillment Trap}

Direct-to-consumer retail can appear attractive because customers pay immediately and the advertised price is visible. But payment fees, shipping, packaging, customer support, discounts, returns, replacements, inventory storage, and advertising can consume the surplus. Paid campaigns may generate impressive sales while the 1st-order contribution remains negative.

Model the entire \emph{order-to-retained-sale} path: qualified traffic, conversion, average order value, refund and return rate, fulfilled-order cost, CAC, repeat purchase, and cash tied up in inventory. A promotional customer who never orders again should not be assigned an invented lifetime value. A better product may increase repeat purchases enough to finance acquisition, but only if that behavior has been measured.

\subsection*{56. SaaS: Recurring Contribution and the Churn Clock}

SaaS can have relatively low marginal software distribution cost, but cloud, support, onboarding, security, AI inference, customer success, and compliance can be material. Model monthly or annual recurring revenue separately from cash prepayments and nonrecurring implementation fees. Track gross revenue retention, account retention, expansion, and contribution margin after cost of service.

CAC payback and lifetime contribution are central when acquisition is front-loaded and value is recurring. The risk is assuming mature retention for a new product. A 10-customer pilot has not established a stable churn hazard. The firm should build survival curves by cohort, keep a conservative cash forecast, and avoid buying growth on the promise of unobserved renewals.

\subsection*{57. Usage-Based AI Systems: Volume Can Increase Costs and Risks}

An AI application may incur a variable API or inference cost every time a customer makes a request. Heavy use can indicate product value yet reduce margins if the pricing metric does not track that cost. Error recovery, human oversight, latency, data protection, and model-provider price changes create further variability.

The useful metric is \textbf{cost per successful, verified task}, not cost per token. A tool that becomes 20 percent cheaper per generation while doubling error-handling time deteriorates overall economics. Analyze gross margin and quality by task type, customer, and usage band. Hybrid commitments, fair-use provisions, and usage caps can protect both sides from surprises if transparently disclosed.

\subsection*{58. Consulting: Utilization and the Scarcity of Expertise}

A consulting firm's key assets are people, judgment, methods, and client relationships. Revenue per consultant and billable utilization matter, but high utilization does not imply healthy economics. Nonbillable proposal writing, training, quality assurance, management, hiring, and recovery are needed to sustain delivery. If every professional is scheduled at full capacity, a single delayed project can cascade into missed commitments.

Track realized fee per project, actual expert hours including sales and revision, direct labor cost, subcontractors, travel, collection timing, project overruns, and client retention. A firm can improve margins through specialization and reusable methodology without claiming its human labor has magically become free. A fixed-fee project should include an explicit change-control mechanism so the buyer and seller know what is included.

\subsection*{59. Marketplaces: Gross Activity Is Not Net Revenue}

A platform may enable transactions worth ₹50 crore and collect a 10 percent commission, but revenue recognition depends on whether it acts as principal or agent and the actual contractual facts. Regardless of accounting presentation, the platform's economics depend on payment processing, acquiring both sides, trust and safety, dispute resolution, fraud, refunds, and liquidity. A large transaction value does not imply a large take-home profit.

A platform may subsidize 1 side to create value for the other. That can be rational if genuine cross-side network effects produce more completed, profitable transactions. It becomes speculative when the founder counts registrations but cannot demonstrate matching, repeat activity, or a path to positive system contribution.

\subsection*{60. Manufacturing: Throughput, Yield, and Capital Intensity}

A manufacturing operation combines raw material, conversion labor, energy, quality control, maintenance, equipment, inventory, and production scheduling. Fixed equipment costs may be large, and yield losses or defect rates can dominate small differences in selling price. Capacity utilization matters, but near-full utilization may increase downtime, expedite costs, and quality failures.

The firm must separate gross profit from investment return. A plant can report respectable operating profit but earn inadequate returns on the large amount of capital tied up in equipment and working capital. Unit economics should eventually connect to return on invested capital and asset replacement, not only cash margin per item.

\subsection*{61. Education: Content Economics Are Not Transformation Economics}

The marginal cost of distributing a recorded explanation can be low, while genuine learning support, feedback, valid assessment, and student success have material human and technical costs. An education venture is therefore neither simply a media business nor simply a classroom rental. Its economic product is a sustained learning system for a defined learner population.

An institution can sell self-paced content, live cohorts, graded practice, private mentoring, curriculum delivery to colleges, and employer-funded training. Each has different unit economics and claims about outcomes. Enrollment counts and completion rates do not establish that learners became capable. A model that improves contribution by abandoning essential feedback may destroy the value proposition that justified the price. \textbf{Cost per verified successful outcome} should supplement, not replace, ordinary financial metrics.

\subsection*{62. Colleges, Enterprise Contracts, and Government Projects}

Institutional programs may involve larger contract values than individual courses, but they often require procurement, scheduling, customization, reporting, compliance, administrative coordination, and receivable financing. A signed memorandum or enthusiastic meeting is not equivalent to an enforceable paid agreement. The buyer may approve a project late, pay only after milestones, or change requirements mid-delivery.

The founder should model \emph{contract-level contribution and cash timing}, including direct labor and travel, scope risk, collection probability, and customer concentration. A ₹10 lakh enterprise program demanding 500 expert-hours and 6 months of sales may be economically weaker than a ₹4 lakh standardized project requiring 80 hours. The comparison is not determined by nominal deal size alone.
